% -----------------------------------------------------------------------------
% This file is automatically generated.
% Do not edit manually.
% Source: notebooks/support/auxiliary-variables-and-stratification/support.Rmd
% -----------------------------------------------------------------------------


\noindent We first create a Classical Variables Sampling ("CVS") object. An object is a container that gathers all the data relevant to a particular sampling application. The \ttblue{FSaudit} package recognizes three such objects. We work with the CVS object \ttblue{cvs\_obj} in this chapter, and with the attribute \ttblue{att\_obj} and MUS \ttblue{mus\_obj} objects in Chapter 5.

\noindent We create the CVS object \ttblue{mySample} by running the \ttblue{cvs\_obj} function, and at the same time fill it with information on the required sample size \ttblue{n}, the relevant book values \ttblue{bv}, identifiers \ttblue{id}, and the random seed number \ttblue{seed}.

\par\addvspace{\topsep}
\begingroup
\fvset{listparameters={%
  \setlength{\topsep}{0pt}%
  \setlength{\partopsep}{0pt}%
  \setlength{\parsep}{0pt}%
  \setlength{\itemsep}{0pt}%
}}
\begin{Verbatim}[breaklines=true,formatcom=\color{ada_blue}]
mySample <- cvs_obj(
  n = 400,
  bv = inventoryData$bv,
  id = inventoryData$item,
  seed = 12345
)
\end{Verbatim}
\endgroup
\par\addvspace{\topsep}

\noindent The CVS object has 19 different attributes attached to it, that will be filled as we proceed from sample planning, through stratification and selection to evaluation. The following is a list of all attributes.
